\documentclass[a4paper,11pt]{article}
\usepackage[margin=27mm]{geometry}
\usepackage{fontspec,amsmath,amssymb,amsthm}
\setmainfont{STIXGeneral}
\usepackage{unicode-math}\setmathfont{STIX Two Math}
\setlength{\emergencystretch}{2em}
\usepackage[hidelinks]{hyperref}
\newcommand{\Sym}{\operatorname{Sym}}\newcommand{\gen}[1]{\langle #1\rangle}
\newtheorem{theorem}{Theorem}\newtheorem{lemma}[theorem]{Lemma}\newtheorem{corollary}[theorem]{Corollary}
\title{Gluing symmetric-group markings and detecting a failed intersection}
\author{Human Compression project}\date{Private application draft; independent review pending}
\begin{document}\maketitle
\begin{abstract}
We derive the principal constructions and uniform counterexample in
Leemans--Mulpas, \emph{Two gluing methods for string C-group representations
of the symmetric groups}, arXiv:2504.17535v1. The common inputs are the
supported-orbit gluing theorem and the exact subgroup formulas for two copies
of a tree. The remaining local arguments are given explicitly, including the
small dihedral boundary and the parity obstruction. The counterexample is
the source's already proved disproof, not a new answer to a conjecture.
The exact source clauses, ancillary exclusions and limited formal coverage
are recorded in the accompanying source map; independent review is pending.
\end{abstract}

\section{Conventions and precise common inputs}
A C-marking is a generating family of distinct nonidentity involutions for
which the subgroups on any two label subsets intersect in the subgroup on
their common labels. It is string when marks at distance at least two in the
chosen order commute. All permutation products use one consistent action
convention; our membership and order arguments do not depend on choosing
left rather than right actions. Conjugation transports the support of a
transposition to its image under the conjugating permutation.

We use the following results with their complete proofs in the two companion
common papers; they are part of the mathematical dependency, not free source
lemmas.
\begin{enumerate}
\item A pure bridge between two disjoint permutation groups generates the
full supported symmetric group on the union of its endpoint orbits. The
complementary induced action is an independent direct factor.
\item Two C-markings glued at one point remain C when every selected subgroup
contains the symmetric group supported on its orbit of that point. A pure
leaf mark guarantees this. Reversing one string order gives the literal
first gluing operation of the source.
\item Two aligned path markings of different lengths give the independent
symmetric groups on the two paths. Equal paths give a diagonal symmetric
group. These statements include all subset intersections.
\item A path of paired vertices with one independent flip at an endpoint gives
the string marking of $C_2\wr S_n$. A path doubled except at its terminal
edge, followed by the full row swap, gives the string marking of
$S_n\wr C_2$. The first pure tail edge on the latter is proved separately
by its complete selected-subgroup formulas; further tail edges use item~2.
\end{enumerate}
We also retain the standard end-parabolic criterion for string markings,
stated as Proposition~2.3 in the source (McMullen--Schulte, Proposition
2E16(a)): the two end-deletion markings must be C and their intersection
must be the subgroup on the middle labels. This implies every subset
intersection. We use it only where those three hypotheses are explicitly
verified below.

The source's Theorem~1.1 and Lemmas~3.1--3.2 now follow from items~1--2,
including their disconnected-component clauses. Its Lemmas~4.2--4.5 follow
from items~3--4. The same common proofs recover Lemma~4.9 and Proposition~4.10,
with the source's actual markings: for $r=3$, the latter wreath marking is
$(12)(56),(23),(15)(26)(34)$, and its first added tail is attached to point~4.
The marked correspondence, rather than the abstract wreath-product order,
is what identifies these applications. The source's own first impure-leaf
counterexample is fully proved in the first companion paper.

\section{The second limitation of leaf attachment}
The source gives a second small example on physical pages~6--7. Its original
marks on $\{1,\ldots,7\}$ are
\[
 a=(12)(56),\qquad b=(23)(45)(67),\qquad c=(34).
\]
They satisfy the string relation, since $a$ and $c$ have disjoint supports.
Both end groups have rank two and are C. To check their intersection, put
$H=\gen{b,c}$. The product $bc$ has order four, so $|H|=8$.
The orbit of the two-element set $\{2,3\}$ under $H$ is exactly
\[
 \bigl\{\{2,3\},\{2,4\},\{3,5\},\{4,5\}\bigr\}.
\]
These four sets are closed under $b,c$ and reached from the first, so its
stabilizer has order two and equals $\gen{b}$. Every element of $H$ fixes~1,
whereas $\gen{a,b}$ preserves $\{1,2,3\}$. Their intersection must therefore
stabilize $\{2,3\}$, and is exactly $\gen{b}$. The retained end-parabolic
criterion proves that the original marking is C.

Its full group is $S_7$. Indeed, starting with $c$ and successively conjugating
by $b,a,b,a,b$ gives the transpositions on the pairs
\[
 \{3,4\},\ \{2,5\},\ \{1,6\},\ \{1,7\},\ \{2,7\},\ \{3,6\}.
\]
These are the edges of a spanning path on the seven points and hence generate
all of $S_7$. Now adjoin point~0 and the mark $x=(01)$ before $a$.
The resulting marking still satisfies the string relations. Since
$(xa)^3=(56)$, we have
\[
 \gen{x,a}=\Sym\{0,1,2\}\times\gen{(56)}.
\]
Its intersection with the original group, which is $S_7$ fixing~0, is
$\gen{(12),(56)}$, of order four. The common-label subgroup $\gen{a}$
has order two. This proves exactly the source's second failed intersection;
no finite computation is a premise. Both small examples express the same
limitation: a distinguished endpoint edge need not be a pure mark.

\section{A folded path followed by a tail}
Let $n\geq2$, take paired vertices $a_1,\ldots,a_n,b_1,\ldots,b_n$,
and a new point $c$. Put
\[
 f=(a_1b_1),\qquad
 d_j=(a_ja_{j+1})(b_jb_{j+1})\ (1\leq j<n),\qquad
 t=(a_nc).
\]

\begin{lemma}\label{lem:folded-tail}
The ordered marking $(f,d_1,\ldots,d_{n-1},t)$ is a string C-marking of
$S_{2n+1}$. Any further pure tail at $c$ preserves the string C-marking and
gives the full symmetric group on the enlarged vertex set.
\end{lemma}
\begin{proof}
Distant generators have disjoint supports, except for doubled path marks
whose distant edges are disjoint on both rows; hence the string relations
hold. Omitting $t$ leaves the paired-path marking of $C_2\wr S_n$, which
is C by the common theorem. Omitting $f$ leaves two aligned paths of lengths
$n$ and $n-1$, one on $A\cup\{c\}$ and one on $B$. Its group is
$\Sym(A\cup\{c\})\times\Sym(B)$ and its marking is C.

An element of the first end group fixes $c$ and preserves the pairs
$\{a_j,b_j\}$. If it also lies in the second end group, it preserves both
rows separately. Thus it can flip no pair and must use the same permutation
of their indices in both rows. Conversely all such diagonal permutations
belong to both end groups and are generated by $d_1,\ldots,d_{n-1}$.
The end intersection is therefore exactly the required middle subgroup.
The retained end-parabolic criterion proves the full intersection property.

The paired-path group is transitive on $A\cup B$ and fixes $c$.
Conjugating the pure transposition $t$ by this group supplies every
$(c\ x)$ with $x\in A\cup B$. These generate $S_{2n+1}$, proving the
full-group claim without a maximal-subgroup input. At $c$ only the pure
last mark acts, so the common leaf theorem proves every further extension.
\end{proof}

Set $n=h+1$. The first graph in source Proposition~4.7 has the folded labels
$h,\ldots,1,0,1,\ldots,h$, followed by a pure path of labels
$h+1,\ldots,r-1$. The lemma gives its initial degree $2h+3$ and rank $h+2$;
adding $r-h-2$ vertices gives rank $r$, degree $r+h+1$ and group $S_{r+h+1}$.
This proves the source's entire stated range $r\geq3$, $1\leq h\leq r-2$.
The harmless $h=0$ base used within its proof is simply the adjacent
transposition marking of $S_3$.

\section{The auxiliary two-component family}
For $r\geq2$, let the first row contain
$a_1,\ldots,a_r,c,d$ and the second contain $b_1,\ldots,b_r$.
Define
\[
 u_j=(a_{j+1}a_{j+2})(b_{j+1}b_{j+2})\quad(0\leq j\leq r-3),
\]
\[
 u_{r-2}=(a_{r-1}a_r)(b_{r-1}b_r)(cd),\qquad
 u_{r-1}=(a_rc).
\]
An empty range of the first display means there are no such generators.
These are precisely the labels of source Lemma~4.6.

\begin{lemma}\label{lem:auxiliary}
This is a string C-marking. Its group is $S_{r+2}\times S_r$ for $r\geq3$,
and the dihedral group of order eight for $r=2$.
\end{lemma}
\begin{proof}
Put $v=(cd)$. For $r\geq3$, cubing $u_{r-3}u_{r-2}$ isolates $v$:
on each of the two long rows the product is a three-cycle, while on $\{c,d\}$
it is $v$. Thus the end group omitting $u_{r-1}$ is
\[
 M_r=\operatorname{diag}S_r\times\gen{v}.
\]
Its two factors act independently. Conjugating the pure bridge
$u_{r-1}=(a_rc)$ by these factors gives all crossing transpositions between
$A$ and $\{c,d\}$. They generate the full supported symmetric group on
$A\cup\{c,d\}$. Cancelling each first-row adjacent transposition from the
corresponding doubled mark then supplies every adjacent transposition on
$B$. This proves the asserted direct-product group without invoking an
external classification of its first-row action.

The string commutations are visible in the displayed supports. For $r=2$
the product of the two involutions is a four-cycle on the first row and a
transposition on the second, hence has order four. The group is dihedral
of order eight and two distinct involutions satisfy the intersection property.
For $r=3$, deleting the first mark gives that rank-two dihedral marking.
Deleting the last gives $M_3$, an endpoint central-involution extension of
the diagonal $S_3$ path. Its two generators have product of order six, so it
is a dihedral group of order twelve and is C.

To check their intersection, the first end group fixes $a_1,b_1$.
Inside $M_3$, the subgroup fixing those points is generated by
\[
 p=(a_2a_3)(b_2b_3),\qquad v=(cd).
\]
The other end group is generated by $q=pv$ and $(a_3c)$.
On the four points $a_2,a_3,c,d$ it preserves the two blocks
$\{a_2,d\}$ and $\{a_3,c\}$ as a set of blocks.
Neither $p$ nor $v$ preserves this block system, while $1$ and $pv=q$ do.
Its intersection with $\gen{p,v}$ is therefore exactly $\gen{q}$, the common
middle subgroup. The end-parabolic criterion proves the rank-three case.

For $r\geq4$, deleting $u_0$ gives the preceding family with parameter
$r-1$, together with the fixed points $a_1,b_1$. It is C by induction.
The group $M_r$ omitting the last mark has a C-marking by the standard
endpoint central-involution extension theorem, stated in source
Proposition~2.2, originally Fernandes--Leemans--Mixer (2012),
Proposition~5.3. This is the only retained extension input in this proof:
the underlying diagonal path marking is C, and only its last generator is
multiplied by the disjoint involution $v$.
Its subgroup fixing $a_1,b_1$ is
$\operatorname{diag}S_{r-1}\times\gen{v}$. The middle generators
$u_1,\ldots,u_{r-2}$ generate exactly that subgroup: their adjacent last pair
again has cube $v$, and their projections give all adjacent transpositions
on indices $2,\ldots,r$. The two end groups therefore intersect in the
middle subgroup, because the first fixes $a_1,b_1$ and the reverse inclusion
is immediate. The retained end-parabolic criterion finishes the induction.
\end{proof}

\section{The uniform counterexample as a parity obstruction}
Let $h\geq1$, $r\geq h+4$, put $n=h+1$ and $m=r-h-2\geq2$.
Use a folded core $A=\{a_1,\ldots,a_n\}$,
$B=\{b_1,\ldots,b_n\}$, a point $c$, and two further rows
$X=\{x_1,\ldots,x_m\}$, $Y=\{y_1,\ldots,y_m\}$.
All these sets are disjoint. Define
\[
 f=(a_1b_1),\qquad d_j=(a_ja_{j+1})(b_jb_{j+1})\ (1\leq j\leq h),
\]
\[
 t=(a_nc),\qquad \tau=\prod_{j=1}^{m}(x_jy_j),\qquad
 s=t\tau,\qquad u=(cx_1),
\]
\[
 e_j=(x_jx_{j+1})(y_jy_{j+1})\quad(1\leq j<m).
\]
In this order the labels are $0,1,\ldots,h,h+1,h+2,h+3,\ldots,r-1$.
Reversing the numbering of $X,Y$ identifies them with the source's displayed
numbers $1,\ldots,2m$ on physical page~11. Thus these are the literal graphs
of source Proposition~4.8 and of its cited graph~(X), in different vertex names.

\begin{theorem}\label{thm:counter}
The displayed marking is a string generating set of $S_{2r-1}$, but is not
a C-marking. More precisely, if
\[
 P=\gen{f,d_1,\ldots,d_h,s},\quad
 Q=\gen{d_1,\ldots,d_h,s,u},\quad
 K=\gen{d_1,\ldots,d_h,s},
\]
then $\tau\in P\cap Q$ and $\tau\notin K$.
\end{theorem}
\begin{proof}
The displayed supports give distinct nonidentity involutions. Distant folded-core marks commute by their
supports, and the full swap $\tau$ commutes with all doubled $e_j$.
The pure bridge $u$ is adjacent only to $s$ and $e_1$. All other distant
pairs have disjoint supports, proving the string relations.

Project $P$ to $A\cup B\cup\{c\}$. Its marked image is the full
$S_{2n+1}$ of Lemma~\ref{lem:folded-tail}. The actual generator $f$ is a pure
transposition, supported entirely on this set. Choose an element of $P$
whose local image sends its endpoints to $a_n,c$. Its conjugate of $f$ is
therefore the actual pure transposition $t$, irrespective of that element's
outside action. Hence $t\in P$, and $\tau=st\in P$.
The same argument provides every transposition on the core, so $P$ contains
the full supported core symmetric group.

For $Q$, consider its component
$U=A\cup\{c,x_1,y_1\}$. Its local marks include the adjacent transpositions
on $A$, the involution $(a_nc)(x_1y_1)$, and the bridge $(cx_1)$.
The cube of the product of the last adjacent transposition on $A$ with
$(a_nc)(x_1y_1)$ is $(x_1y_1)$, since the other factor is a three-cycle.
Thus the local image contains both $\Sym(A\cup\{c\})$ and the independent
$\Sym\{x_1,y_1\}$. The pure bridge $(cx_1)$ joins them, giving all of
$\Sym(U)$. Since $u$ is an actual pure transposition, conjugating it by a
suitable element of $Q$ gives the actual $t=(a_nc)$. Therefore
$\tau=st\in Q$ as well.

Every generator $d_j$ is even on the core $A\cup B\cup\{c\}$ and fixes
$X\cup Y$. The generator $s$ restricts to the odd transposition $t$ on the
core and to $\tau$ on $X\cup Y$. In every word on the generators of $K$,
the exponent of $\tau$ is consequently the sign of its core restriction.
If that restriction is the identity, the exponent must be even. The nontrivial
permutation $\tau$ alone has identity core restriction, so it cannot belong
to $K$. This proves the failed intersection on the two explicitly specified
label subsets.

Finally, the full group contains the supported symmetric group on the core,
as proved above, and the group generated on $X\cup Y$ by $\tau$ and all
$e_j$. The latter is transitive: the doubled path moves within either row,
and $\tau$ interchanges the rows. These two supported groups have disjoint
supports and the pure bridge $u$ connects them. The common pure-bridge lemma
gives the full symmetric group on their union. Its degree is
$2n+1+2m=2r-1$, proving the asserted whole group.
\end{proof}

This is the source's already established uniform counterexample to its
Conjecture~4.1 and to the indicated earlier Table~6 conjecture. The proof
isolates the reusable mechanism: the smaller selected group ties one outside
involution to a character of its local action, whereas each larger selected
group separates that involution by an actual supported transposition.
No new conjecture answer is claimed. The source's concluding numerical
stability sequence and bibliographic preliminaries are not deductions of this
application; they remain visible exclusions in the source map.
\end{document}
